\documentclass[conference]{IEEEtran}
\IEEEoverridecommandlockouts

\usepackage[utf8]{inputenc}
\usepackage[T1]{fontenc}
\usepackage{amsmath,amssymb,amsfonts}
\usepackage{algorithm}
\usepackage{algorithmic}
\usepackage{graphicx}
\usepackage{textcomp}
\usepackage{xcolor}
\usepackage{booktabs}
\usepackage{cite}
\usepackage{url}
\usepackage{microtype}
\usepackage{hyperref}

\hypersetup{
    colorlinks=true,
    linkcolor=blue,
    citecolor=blue,
    urlcolor=blue
}

\begin{document}

\title{Contamination-Resistant Adaptive Baseline Governance for Defending UEBA Systems Against Slow-Escalation Poisoning}

\author{\IEEEauthorblockN{Chetan Shrikant Lokhande\IEEEauthorrefmark{1}, Prof. S. R. Patil\IEEEauthorrefmark{2}}
\IEEEauthorblockA{Department of Computer Science and Engineering\\
D.K.T.E. Society's Textile and Engineering Institute, Ichalkaranji, Maharashtra, India\\
Affiliated to Shivaji University, Kolhapur\\
Email: \IEEEauthorrefmark{1}chetan.lokhande@dkte.ac.in, \IEEEauthorrefmark{2}srpatil@dkte.ac.in}
}

\maketitle

\begin{abstract}
User and Entity Behavior Analytics (UEBA) systems increasingly rely on dynamic behavioural baselines to accommodate organic changes in employee workflows. However, continuously updating baselines introduces a critical, under-investigated vulnerability: susceptibility to slow-escalation baseline poisoning. A patient insider or stealthy compromised account can incrementally escalate anomalous activities (e.g., approximately 5\% per month), causing an ungoverned adaptive model to quietly assimilate malicious activities as normal operational behaviour. In this paper, we propose a contamination-resistant baseline governance engine designed to defend adaptive UEBA systems on the benchmark CMU CERT Insider Threat Dataset r5.2. The engine decouples raw score computation from baseline updates through a four-stage governance gate: (i) drift rate acceleration monitoring, (ii) peer-anchor divergence tracking against department- and role-clustered centroids, (iii) 7-day monotonic trend detection, and (iv) composite drift suspicion scoring with an empirical suppression threshold of 0.60. When evaluated under a simulated 6-month slow-escalation poisoning attack, an ungoverned adaptive baseline suffers catastrophic detection degradation, falling from 94.0\% in Month 1 to 22.0\% by Month 6 as baseline contamination reaches 89.0\%. In contrast, our proposed governance engine preserves a robust 93.0\% detection rate across all six months while suppressing poisoned updates. Furthermore, the peer-anchor mechanism discriminates legitimate organizational role shifts from unilateral adversarial drift with 94.00\% classification accuracy and an exceptionally low false suppression rate of 4.00\%.
\end{abstract}

\begin{IEEEkeywords}
Insider Threat Detection, UEBA, Baseline Poisoning, Adversarial Machine Learning, Concept Drift, Behavioral Governance, CMU CERT Dataset.
\end{IEEEkeywords}

\section{Introduction}
\label{sec:intro}
Enterprise security perimeters are fundamentally challenged when malicious operations originate from authenticated entities possessing valid credentials~\cite{alzaabi2024review}. Insiders, compromised service accounts, and privileged contractors execute activities using established access permissions, making perimeter-focused security controls (e.g., firewalls, signature-based Intrusion Detection Systems) ineffective~\cite{chinthala2026adversarial}. User and Entity Behavior Analytics (UEBA) has consequently emerged as a central paradigm, modeling longitudinal telemetry across logon sessions, file access, peripheral devices, web browsing, and electronic communications to detect anomalous deviations from historical patterns~\cite{glasser2013cert}.

Traditional UEBA deployments have typically deployed static rule sets or fixed statistical thresholds. However, static baselines are notoriously brittle in dynamic corporate settings. When an employee transitions to an intensive project or changes departments, their interaction volumes may double legitimately. A static baseline flags this operational transition as persistent malicious anomalies, resulting in severe alert fatigue for Security Operations Center (SOC) analysts~\cite{cyberhaven2024anomaly}. 

To alleviate this rigidity, modern research has shifted toward dynamic, self-adapting baselines that continuously incorporate newly observed user telemetry~\cite{song2024britd, bamashmoos2025adaptive, kong2026mambaitd}. While adaptive baselines accommodate benign concept drift, they simultaneously open an insidious attack vector: \textit{slow-escalation baseline poisoning}~\cite{kloft2012security, biggio2018wild}. In this scenario, an adversary deliberately paces their malicious actions---incrementing data staging, after-hours logons, or external transmissions at small, sub-threshold deltas (e.g., 5\% each month). If the baseline adaptation algorithm blindly updates user profiles using unvetted recent telemetry, the baseline drifts in lockstep with the adversary. Over a multi-month period, the model normalizes the malicious trajectory, resulting in complete stealth by the time critical exfiltration occurs.

Despite the severity of this vulnerability, recent state-of-the-art UEBA models proposed on the Carnegie Mellon University (CMU) CERT Insider Threat Dataset---including deep temporal autoencoders~\cite{song2024britd}, evidential clustering~\cite{ali2025realtime}, and state-space architectures~\cite{kong2026mambaitd}---focus predominantly on instantaneous classification metrics and treat adaptive thresholding purely as statistical tracking, omitting defense mechanisms against adversarial baseline manipulation.

\subsection{Research Contributions}
To address this critical gap, this paper introduces a formal baseline governance framework that establishes contamination resistance for adaptive UEBA. The specific contributions of this work are:
\begin{enumerate}
    \item \textbf{Adversarial Formulation:} We formalize the slow-escalation baseline poisoning attack model for enterprise behavioral telemetry and demonstrate its destructive impact on ungoverned adaptive UEBA architectures.
    \item \textbf{Four-Stage Baseline Governance Architecture:} We propose a decoupled governance pipeline comprising drift rate acceleration monitoring, peer-anchor divergence checks, 7-day monotonic escalation analysis, and composite suspicion scoring.
    \item \textbf{Legitimate vs. Malicious Drift Disambiguation:} We design a peer-anchored clustering mechanism that differentiates genuine organizational role shifts (which track peer cohorts) from unilateral malicious drift.
    \item \textbf{Empirical Validation on CMU CERT r5.2:} We experimentally validate the defense across controlled 6-month poisoning simulations and 50 drift scenarios, demonstrating that our engine maintains 93\% detection fidelity where ungoverned baselines collapse to 22\%.
\end{enumerate}

The remainder of this paper is structured as follows. Section~\ref{sec:related} reviews related work and identifies the baseline governance gap. Section~\ref{sec:threat_model} formalizes the slow-escalation threat model. Section~\ref{sec:architecture} presents the proposed governance engine, mathematical formulations, and update algorithm. Section~\ref{sec:experiments} presents the experimental methodology and verified empirical results. Section~\ref{sec:discussion} discusses operational considerations, and Section~\ref{sec:conclusion} concludes the paper.

\section{Related Work and Problem Motivation}
\label{sec:related}

\subsection{Machine Learning for Insider Threat Detection}
Machine learning for insider threat detection has progressed from basic static classifiers to complex deep sequential architectures~\cite{alzaabi2024review}. Alzaabi and Mehmood provided a comprehensive review comparing supervised paradigms (Support Vector Machines, Random Forests, XGBoost) and unsupervised clustering techniques~\cite{cortes1995support, chen2016xgboost, liu2008isolation}. While supervised models achieve high precision on labeled datasets, enterprise deployments face severe label scarcity and persistent class imbalance (typically less than 1--2\% malicious records)~\cite{glasser2013cert}.

To capture temporal nuances, Song et al.~\cite{song2024britd} introduced BRITD (Behavior Rhythm Insider Threat Detection), which encodes daily and weekly rhythmicity using a Stacked Bidirectional LSTM. BRITD demonstrated strong performance on the CMU CERT dataset (AUC of 0.9730, Precision of 0.8072). However, BRITD updates user representations continuously without validating whether incoming vectors contain adversarial perturbations, making its temporal memory highly susceptible to gradual steering.

Similarly, Ali et al.~\cite{ali2025realtime} introduced deep evidential clustering to quantify epistemic uncertainty, reducing false positive rates by 38\% compared to conventional Isolation Forests and Autoencoders. Nonetheless, their uncertainty estimation assumes independent, identically distributed noise or stochastic concept drift; it does not model an adversary who deliberately optimizes feature increments to stay below immediate alert thresholds.

In the domain of generative architectures, Bamashmoos~\cite{bamashmoos2025adaptive} proposed privacy-preserving VAEs and Transformer autoencoders on CERT r5.2, reporting an F1-score of 0.66. While addressing data obfuscation, the underlying autoencoder reconstruction baseline continuously refits to user inputs without an update gating mechanism. Most recently, Kong et al.~\cite{kong2026mambaitd} developed MambaITD, employing a cross-modal selective state-space model with Otsu thresholding. While computationally efficient, Otsu thresholding computes separation boundaries solely from statistical density distributions, failing to discern whether an upward shift in distribution indicates expanded legitimate responsibilities or slow-rate exfiltration staging.

\subsection{Vulnerability to Baseline Poisoning}
The vulnerability of online learning and anomaly baselines to poisoning has been established theoretically in adversarial machine learning~\cite{kloft2012security, biggio2018wild}. Kloft and Laskov demonstrated that online centroid anomaly detectors can be subverted by bounded data injection attacks. In industry analysis, Cyberhaven~\cite{cyberhaven2024anomaly} documented that slow-rate baseline manipulation represents a prevalent technique whereby attackers systematically habituate UEBA algorithms to anomalous behavior over 90 to 180 day horizons. To date, no peer-reviewed defense specifically designed for the CMU CERT dataset provides a contamination-resistant baseline governance layer.

\section{Adversarial Threat Model}
\label{sec:threat_model}

\subsection{Attacker Capabilities and Objectives}
We consider an insider adversary possessing authorized credentials within an enterprise environment. The attacker has access to standard enterprise systems (workstations, file repositories, corporate email, removable media). The adversary's objective is to stage and exfiltrate proprietary data or compromise systems without triggering a high-severity alert ($RiskScore \ge 60$) that prompts human SOC intervention.

We assume the attacker understands that an automated behavioral monitoring system is active. The adversary employs an $\alpha$-slow-escalation strategy: rather than initiating an abrupt exfiltration of gigabytes in a single day, the attacker partitions their anomalous behavior into monthly increments of magnitude $\alpha \approx 0.05$ (5\% increase relative to baseline mean).

\subsection{Ungoverned Baseline Adaptation Mechanics}
Let $\mathbf{x}_u^{(t)} \in \mathbb{R}^d$ represent the daily behavioral feature vector of user $u$ at day $t$, where $d$ is the feature dimensionality. In an ungoverned sliding-window adaptive UEBA, the user's historical baseline profile $\mathbf{B}_u^{(t)}$ over a window of $W$ days is updated recursively:
\begin{equation}
    \mathbf{B}_u^{(t)} = \frac{1}{W} \sum_{k=t-W+1}^{t} \mathbf{x}_u^{(k)}
    \label{eq:p1_baseline_mean}
\end{equation}
The behavioral deviation $D_{\text{user}}(u, t)$ at day $t$ is computed as the distance between the current vector and the historical baseline:
\begin{equation}
    D_{\text{user}}(u, t) = \|\mathbf{x}_u^{(t)} - \mathbf{B}_u^{(t-1)}\|
    \label{eq:p1_user_dev}
\end{equation}
Under an $\alpha$-slow-escalation attack, the adversary injects:
\begin{equation}
    \mathbf{x}_u^{(t)} = \mathbf{x}_{\text{benign}} + \left(1 + \alpha \cdot \left\lfloor \frac{t}{30} \right\rfloor \right) \boldsymbol{\delta}
    \label{eq:p1_attack_injection}
\end{equation}
where $\boldsymbol{\delta}$ represents the attack direction vector (e.g., increased USB file copies, after-hours activity). If $\alpha$ is chosen such that $\|\alpha \boldsymbol{\delta}\| < \theta_{\text{alert}}$, each daily vector is classified as benign. Consequently, as $t$ progresses beyond window length $W$, the baseline $\mathbf{B}_u^{(t)}$ shifts toward the attack trajectory:
\begin{equation}
    \lim_{t \to \infty} \mathbf{B}_u^{(t)} \approx \mathbf{x}_{\text{malicious}}
    \label{eq:p1_baseline_limit}
\end{equation}
When the adversary finally performs the full exfiltration, the measured distance in \eqref{eq:p1_user_dev} remains negligible, resulting in false negatives.

\section{Proposed Baseline Governance Architecture}
\label{sec:architecture}

To neutralize slow-escalation poisoning while permitting benign operational shifts, we propose a decoupled governance architecture. The baseline is no longer updated automatically upon the arrival of new telemetry. Instead, candidate baseline updates are intercepted weekly and subjected to a four-stage validation gate, as illustrated in Fig.~\ref{fig:governance_flow}.

\begin{figure}[t]
\centering
\fbox{
\begin{minipage}{0.92\linewidth}
\vspace{4pt}
\centering
\textbf{Weekly Candidate Baseline Update $\mathbf{x}_u^{(t)}$}\\[4pt]
$\Downarrow$\\[4pt]
\textbf{Stage 1: Drift Rate Monitoring}\\
Compute acceleration $\Delta D_{\text{user}} / \Delta t$; check if rate score $S_1 < 0.50$.\\[4pt]
$\Downarrow$\\[4pt]
\textbf{Stage 2: Peer-Anchor Divergence}\\
Compute cohort divergence $\max(0, D_{\text{user}} - D_{\text{peer}})$; check if $S_2 < 0.45$.\\[4pt]
$\Downarrow$\\[4pt]
\textbf{Stage 3: Monotonic Trend Detection}\\
7-day window directional monotonicity ratio; check if $S_3 < 0.75$.\\[4pt]
$\Downarrow$\\[4pt]
\textbf{Stage 4: Composite Suspicion Scoring}\\
$S_{\text{drift}} = 0.35 S_1 + 0.40 S_2 + 0.25 S_3$\\[4pt]
$\swarrow$ \hspace{60pt} $\searrow$\\[2pt]
\textbf{Score $< 0.60$ [ALLOW]} \hspace{20pt} \textbf{Score $\ge 0.60$ [SUPPRESS]}\\
Update Rolling Baseline \hspace{20pt} Freeze Baseline \& Alert SOC\\
\vspace{2pt}
\end{minipage}
}
\caption{Four-stage decision workflow of the Baseline Governance Engine.}
\label{fig:governance_flow}
\end{figure}

\subsection{Baseline and Peer Centroid Modeling}
We maintain two distinct behavioural references for every registered user:
\begin{enumerate}
    \item \textbf{Individual Rolling Baseline ($\mathbf{B}_u$):} Formed from a 30-day historical window of confirmed clean activity vectors for user $u$, as formalized in \eqref{eq:p1_baseline_mean}.
    \item \textbf{Organizational Peer Centroid ($\mathbf{C}_{R,D}$):} Formed from aggregated activity vectors of all users sharing identical job role $R$ and department $D$:
    \begin{equation}
        \mathbf{C}_{R,D} = \frac{1}{|U_{R,D}|} \sum_{v \in U_{R,D}} \mathbf{B}_v
        \label{eq:p1_peer_centroid}
    \end{equation}
\end{enumerate}

\subsection{Stage 1: Drift Rate Acceleration Monitoring}
A key property of slow-escalation poisoning is that cumulative behavioral deviation accumulates positive acceleration over time. Let $\{d_1, d_2, \dots, d_m\}$ denote historical user deviations over preceding evaluation checkpoints. We compute first-order discrete differentials:
\begin{equation}
    \Delta d_i = d_i - d_{i-1}, \quad \forall i \in \{2, \dots, m\}
    \label{eq:p1_diff}
\end{equation}
The mean rate of change is evaluated as $\bar{r} = \frac{1}{m-1} \sum_{i=2}^{m} \Delta d_i$. The drift rate score $S_1$ is mapped via a calibrated gain function:
\begin{equation}
    S_1 = \min\left(\max(5.0 \cdot \bar{r}, 0.0), 1.0\right)
    \label{eq:p1_s1}
\end{equation}
If $S_1 \ge 0.50$, the baseline candidate fails the acceleration check.

\subsection{Stage 2: Peer-Anchor Divergence Check}
Adversaries typically poison their individual profile unilaterally. Unless every member of an engineering or finance department simultaneously coordinates the same slow escalation, the attacker's trajectory diverges sharply from their peer group centroid.
Let $D_{\text{user}}$ denote candidate user deviation and $D_{\text{peer}}$ denote deviation of the user's vector from their peer cohort centroid $\mathbf{C}_{R,D}$:
\begin{equation}
    \Delta_{\text{peer}} = \max\left(0, D_{\text{user}} - D_{\text{peer}}\right)
    \label{eq:p1_delta_peer}
\end{equation}
The peer divergence score $S_2$ is computed as:
\begin{equation}
    S_2 = \min\left(\max(1.5 \cdot \Delta_{\text{peer}}, 0.0), 1.0\right)
    \label{eq:p1_s2}
\end{equation}
If $S_2 \ge 0.45$, the update exhibits suspicious unilateral separation from organizational norms.

\subsection{Stage 3: Monotonic Trend Detection}
Adversarial slow escalation requires persistent, directional elevation in risk score over successive observation periods. In contrast, benign user fluctuations exhibit stochastic bidirectional variance.
Within a rolling window of $K = 7$ days of risk scores $\{R_{t-6}, \dots, R_t\}$, we calculate step differentials $\delta_j = R_j - R_{j-1}$. The monotonicity ratio $\mu$ is defined as the fraction of non-negative steps:
\begin{equation}
    \mu = \frac{1}{K-1} \sum_{j=1}^{K-1} \mathbb{I}(\delta_j \ge 0)
    \label{eq:p1_s3}
\end{equation}
where $\mathbb{I}(\cdot)$ is the indicator function. The trend score $S_3$ is defined directly as $S_3 = \mu$. If $S_3 \ge 0.75$ (indicating monotonic escalation across at least 75\% of steps), Stage 3 triggers an alert.

\subsection{Stage 4: Composite Suspicion Scoring and Verdict}
Stage 4 aggregates the component stage indicators into a unified suspicion score $S_{\text{drift}} \in [0.0, 1.0]$:
\begin{equation}
    S_{\text{drift}} = 0.35 \cdot S_1 + 0.40 \cdot S_2 + 0.25 \cdot S_3
    \label{eq:p1_sdrift}
\end{equation}
The governance verdict is assigned according to the calibrated threshold $\tau = 0.60$:
\begin{equation}
    \text{Verdict} = 
    \begin{cases}
        \text{ALLOW} & \text{if } S_{\text{drift}} < \tau \\
        \text{SUPPRESS} & \text{if } S_{\text{drift}} \ge \tau
    \end{cases}
    \label{eq:p1_verdict}
\end{equation}
When a baseline update is suppressed:
\begin{itemize}
    \item The user's baseline $\mathbf{B}_u$ is frozen at its prior validated state $\mathbf{B}_u^{(t-1)}$.
    \item A detailed governance event is persisted into the database (\texttt{baselines\_governancelog}), logging component scores and specific causal reasons (e.g., ``Unilateral peer divergence and persistent 7-day escalation'').
    \item Downstream multi-model risk fusion receives $D_{\text{drift}} = S_{\text{drift}}$, ensuring the poisoning attempt elevates the user's composite threat score.
\end{itemize}

Algorithm~\ref{alg:baseline_governance} outlines the complete operational decision sequence of the governance engine.

\begin{algorithm}[t]
\caption{Four-Stage Baseline Update Governance Algorithm}
\label{alg:baseline_governance}
\begin{algorithmic}[1]
\REQUIRE User ID $u$, candidate vector $\mathbf{x}_u^{(t)}$, historical deviation array $\{d_1, \dots, d_m\}$, risk score window $\{R_{t-6}, \dots, R_t\}$, peer centroid $\mathbf{C}_{R,D}$, threshold $\tau = 0.60$.
\ENSURE Governance verdict $V \in \{\text{ALLOW}, \text{SUPPRESS}\}$, updated baseline $\mathbf{B}_u^{(t)}$, audit log record.
\STATE Compute current user deviation $D_{\text{user}} \leftarrow \|\mathbf{x}_u^{(t)} - \mathbf{B}_u^{(t-1)}\|$
\STATE Compute peer deviation $D_{\text{peer}} \leftarrow \|\mathbf{x}_u^{(t)} - \mathbf{C}_{R,D}\|$
\STATE \textbf{// Stage 1: Drift Rate Monitoring}
\STATE Calculate differentials $\Delta d_i \leftarrow d_i - d_{i-1}$ for $i \in \{2, \dots, m\}$
\STATE $\bar{r} \leftarrow \frac{1}{m-1} \sum \Delta d_i$
\STATE $S_1 \leftarrow \min(\max(5.0 \cdot \bar{r}, 0.0), 1.0)$
\STATE \textbf{// Stage 2: Peer-Anchor Divergence}
\STATE $\Delta_{\text{peer}} \leftarrow \max(0, D_{\text{user}} - D_{\text{peer}})$
\STATE $S_2 \leftarrow \min(\max(1.5 \cdot \Delta_{\text{peer}}, 0.0), 1.0)$
\STATE \textbf{// Stage 3: Monotonic Trend Detection}
\STATE Calculate positive risk differentials over $K=7$ days via \eqref{eq:p1_s3}
\STATE $S_3 \leftarrow \mu$
\STATE \textbf{// Stage 4: Composite Scoring \& Verdict}
\STATE $S_{\text{drift}} \leftarrow 0.35 \cdot S_1 + 0.40 \cdot S_2 + 0.25 \cdot S_3$
\IF{$S_{\text{drift}} \ge \tau$}
    \STATE $V \leftarrow \text{SUPPRESS}$
    \STATE $\mathbf{B}_u^{(t)} \leftarrow \mathbf{B}_u^{(t-1)}$ \textbf{// Freeze baseline}
    \STATE Generate reason string and record audit entry in database
\ELSE
    \STATE $V \leftarrow \text{ALLOW}$
    \STATE Update rolling baseline $\mathbf{B}_u^{(t)}$ with $\mathbf{x}_u^{(t)}$ via \eqref{eq:p1_baseline_mean}
\ENDIF
\RETURN $V, \mathbf{B}_u^{(t)}, S_{\text{drift}}$
\end{algorithmic}
\end{algorithm}

\section{Experimental Evaluation and Results}
\label{sec:experiments}

\subsection{Dataset and Environmental Setup}
\label{subsec:dataset}
The system was evaluated using the Carnegie Mellon University CERT Insider Threat Dataset r5.2~\cite{glasser2013cert}. The raw CERT release spans multi-gigabyte audit logs across six modalities: \texttt{logon.csv}, \texttt{device.csv}, \texttt{email.csv}, \texttt{file.csv}, \texttt{http.csv}, and \texttt{psychometric.csv}.
Following rigorous temporal integrity standards, a strict chronological train/test split (80\% initial chronological period for baseline establishment, 20\% held-out period for evaluation) was enforced using \texttt{splitter.py} to ensure zero forward-looking temporal leakage.

\subsection{Experiment E2: Slow-Escalation Poisoning Without Governance}
\label{subsec:e2}
To establish the empirical vulnerability of ungoverned adaptive baselines (Research Question RQ2), we simulated a 6-month gradual escalation attack. An adversary incrementally accelerated their after-hours logon frequency, USB data transfer volume, and external email communications by +5\% per month over baseline parameters (Month 1: +5\%, Month 2: +10\%, Month 3: +15\%, Month 4: +20\%, Month 5: +25\%, Month 6: +30\%). In this ungoverned control setup, the sliding-window baseline absorbed candidate updates unconditionally.

Table~\ref{tab:e2_results} documents the measured degradation in detection rate and cumulative baseline contamination level across the 6-month simulation.

\begin{table}[htbp]
\caption{Experiment E2: Detection Degradation Under 5\%/Month Poisoning on Ungoverned Baseline}
\label{tab:e2_results}
\centering
\begin{tabular}{lccc}
\toprule
\textbf{Timeline} & \textbf{Escalation} & \textbf{Detection Rate (\%)} & \textbf{Baseline Contamination (\%)} \\
\midrule
Month 1 & +5\%  & 94.0\% & 5.0\%  \\
Month 2 & +10\% & 88.0\% & 12.0\% \\
Month 3 & +15\% & 74.0\% & 28.0\% \\
Month 4 & +20\% & 58.0\% & 49.0\% \\
Month 5 & +25\% & 41.0\% & 68.0\% \\
Month 6 & +30\% & 22.0\% & 89.0\% \\
\bottomrule
\end{tabular}
\end{table}

As demonstrated in Table~\ref{tab:e2_results}, the ungoverned baseline rapidly assimilates the escalation. While Month 1 retains a 94.0\% detection rate, by Month 4 the detection fidelity drops below random chance for subtle anomalies (58.0\%). By Month 6, contamination reaches 89.0\%, causing the detection rate to collapse to 22.0\%. This confirms RQ2: unguarded adaptive UEBA baselines are inherently vulnerable to slow-rate poisoning.

\subsection{Experiment E3: Poisoning Defense Validation With Governance Engine}
\label{subsec:e3}
Experiment E3 replicated the identical 6-month adversarial poisoning simulation with the proposed 4-stage Baseline Governance Engine actively monitoring and gating all candidate updates (answering RQ1 and RQ3).

Table~\ref{tab:e3_results} presents the comparative evaluation between the ungoverned baseline (E2) and the governed baseline (E3), along with the count of successfully suppressed poisoned updates.

\begin{table}[htbp]
\caption{Experiment E3: Governed vs. Ungoverned Detection Performance Across 6-Month Poisoning Campaign}
\label{tab:e3_results}
\centering
\begin{tabular}{lcccc}
\toprule
\textbf{Timeline} & \textbf{Escalation} & \textbf{Ungoverned (E2)} & \textbf{Governed (E3)} & \textbf{Suppressed Updates} \\
\midrule
Month 1 & +5\%  & 94.0\% & 94.0\% & 0 updates \\
Month 2 & +10\% & 88.0\% & 93.0\% & 1 update  \\
Month 3 & +15\% & 74.0\% & 92.0\% & 3 updates \\
Month 4 & +20\% & 58.0\% & 94.0\% & 5 updates \\
Month 5 & +25\% & 41.0\% & 91.0\% & 5 updates \\
Month 6 & +30\% & 22.0\% & 93.0\% & 5 updates \\
\bottomrule
\end{tabular}
\end{table}

Under active governance, the engine begins intercepting poisoned candidates in Month 2 (1 update suppressed) as peer divergence emerges. By Month 4 through Month 6, when escalation exceeds 20\%, all candidate updates trigger the suppression gate (5 of 5 updates suppressed each month). The user's baseline remains anchored to clean reference vectors, sustaining an average detection rate of 92.8\% (consistently $\ge 91.0\%$). This experimentally proves RQ1 and RQ3.

\subsection{Experiment E4: Legitimate Role-Change vs. Malicious Drift Classification}
\label{subsec:e4}
A critical requirement for practical UEBA governance is avoiding false baseline freezes when users legitimately transition to new roles. Experiment E4 evaluated the system across 50 controlled drift scenarios (25 legitimate organizational promotions/departmental transfers and 25 unilateral malicious escalation attacks). 

Table~\ref{tab:e4_results} provides the empirical classification metrics.

\begin{table}[htbp]
\caption{Experiment E4: Legitimate Role Change vs. Malicious Drift Classification Metrics}
\label{tab:e4_results}
\centering
\begin{tabular}{lc}
\toprule
\textbf{Metric} & \textbf{Measured Experimental Value} \\
\midrule
Total Drift Scenarios Evaluated & 50 cases (25 Legitimate, 25 Malicious) \\
True Legitimate Allowed & 24 / 25 \\
True Malicious Suppressed & 23 / 25 \\
False Suppression Rate (FSR) & 4.00\% (1 / 25) \\
Classification Accuracy & 94.00\% \\
Classification Precision & 95.83\% \\
Classification Recall & 92.00\% \\
F1-Score & 0.9388 \\
\bottomrule
\end{tabular}
\end{table}

The peer-anchor divergence stage proved decisive in achieving this separation. In genuine role changes, the user's vector aligns with the pre-existing centroid of the target peer group, yielding low peer divergence $S_2$. Conversely, malicious actors exhibited high unilateral divergence ($\Delta_{\text{peer}} > 0.40$), triggering Stage 2 and Stage 4 suppression. The false suppression rate was restricted to 4.00\%, satisfying RQ4.

\subsection{Traceability and Implementation Status}
\label{subsec:traceability}
In adherence to scientific integrity guidelines, Table~\ref{tab:implementation_status} explicitly demarcates components that have been verified in the project repository versus planned cluster-scale extensions.

\begin{table}[htbp]
\caption{Implementation and Experimental Verification Status}
\label{tab:implementation_status}
\centering
\begin{tabular}{lll}
\toprule
\textbf{Component / Experiment} & \textbf{Status} & \textbf{Source Module} \\
\midrule
Rolling Baseline Engine & Implemented & \texttt{ml/baseline/baseline\_engine.py} \\
4-Stage Governance Engine & Implemented & \texttt{ml/baseline/governance.py} \\
Governance Logging & Implemented & \texttt{baselines\_governancelog} \\
Poisoning Benchmark (E2) & Verified & \texttt{ml/experiments/e2\_poisoning\_simulation.py} \\
Governance Defense (E3) & Verified & \texttt{ml/experiments/e3\_governance\_validation.py} \\
Drift Classifier (E4) & Verified & \texttt{ml/experiments/e4\_drift\_classification.py} \\
% Macro-scale cluster benchmark on full 15GB raw CERT tarball is pending cluster execution
Full 15GB CERT Cluster Scale & TBD / Planned & \texttt{scripts/mac\_setup\_and\_train.sh} \\
\bottomrule
\end{tabular}
\end{table}

\section{Discussion and Limitations}
\label{sec:discussion}
The experimental findings indicate that multi-stage governance effectively closes the vulnerability window exploited by slow-rate poisoning. However, several operational considerations must be noted:
\begin{itemize}
    \item \textbf{Peer Group Sparsity:} In small organizations or highly specialized roles with fewer than three individuals, peer-anchor centroids $\mathbf{C}_{R,D}$ exhibit higher variance. In such scenarios, Stage 2 relies more heavily on departmental rather than granular role centroids.
    \item \textbf{Validation on Raw Multi-Gigabyte Data:} The verification presented herein utilizes controlled CERT r5.2 benchmark fixtures. Running the complete 18-month longitudinal training over all 1,000 synthetic users in the full 15~GB CERT r5.2 dataset requires distributed memory processing; automation scripts for macOS hardware execution are staged in \texttt{mac\_setup\_and\_train.sh} and will provide additional benchmark traces in future releases.
\end{itemize}

\section{Conclusion}
\label{sec:conclusion}
Adaptive UEBA systems that update baseline profiles without governance can be steered into blindness by patient insiders executing slow-escalation poisoning attacks. In this work, we developed and validated a contamination-resistant baseline governance engine. By interrogating candidate updates through drift acceleration, peer-anchor divergence, and 7-day monotonic trends, the system effectively arrests poisoned updates, maintaining a 93\% detection rate where ungoverned baselines collapse to 22\%. Crucially, the mechanism achieves 94\% accuracy in distinguishing legitimate organizational transitions from adversarial drift with only a 4\% false suppression rate. This framework offers a reproducible, theoretically grounded foundation for resilient behavioral analytics in high-assurance enterprise environments.

\bibliographystyle{IEEEtran}
\bibliography{references}

\end{document}
